\documentclass{article}
\input{../../../lib.sty}

\title{\texttt{fn ensure\_open\_unit\_interval}}
\author{Michael Shoemate}

\begin{document}
\maketitle

Proves soundness of \texttt{fn ensure\_open\_unit\_interval} in \texttt{mod.rs}.

\section{Hoare Triple}
\subsection*{Precondition}
None.

\subsection*{Pseudocode}
\lstinputlisting[language=Python,firstline=2,escapechar=|]{./pseudocode/ensure_open_unit_interval.py}

\subsection*{Postcondition}
\begin{theorem}
\label{postcondition}
For any input \texttt{x}, \texttt{ensure\_open\_unit\_interval(x)} either returns \texttt{Err(e)},
or returns \texttt{Ok(())} and \texttt{x} lies in the open unit interval $(0,1)$.
\end{theorem}

\section{Proof}
Assume an arbitrary input \texttt{x}.

\begin{lemma}
\label{err-e}
\texttt{ensure\_open\_unit\_interval} returns \texttt{Err(e)} exactly when the guard on
line \ref{line:guard} is true.
\end{lemma}

\begin{proof}
The guard is the only branch in the function. If it is true, the function executes the explicit
\texttt{raise}. If it is false, the function reaches the final \texttt{return ()}.
\end{proof}

\begin{lemma}
\label{ok-out}
If \texttt{ensure\_open\_unit\_interval(x)} returns \texttt{Ok(())}, then \texttt{x} lies in
$(0,1)$.
\end{lemma}

\begin{proof}
If the function returns \texttt{Ok(())}, then by Lemma~\ref{err-e} the guard on
line \ref{line:guard} is false. Therefore \texttt{x} is neither $\le 0$ nor $\ge 1$, so
\texttt{x} lies in $(0,1)$.
\end{proof}

\begin{proof}[Proof of Theorem~\ref{postcondition}]
Holds by Lemma~\ref{err-e} and Lemma~\ref{ok-out}.
\end{proof}

\end{document}
